\documentclass[11pt,letter,notitlepage]{article}
\usepackage[left=2cm, right=2cm, lines=45, top=0.8in, bottom=0.7in]{geometry}
\usepackage{fancyhdr}
\usepackage{fancybox}
\usepackage{graphicx}
\usepackage{pdfpages}
\usepackage{enumitem}
\usepackage{algorithm}
\usepackage{algorithmic}
\renewcommand{\headrulewidth}{1.5pt}
\renewcommand{\footrulewidth}{1.5pt}
\newcommand\Loadedframemethod{TikZ}
\usepackage[framemethod=\Loadedframemethod]{mdframed}
\usepackage{amssymb,amsmath}
\usepackage{mathtools}
\usepackage{amsthm}
\usepackage{thmtools}
\usepackage{bm}
\usepackage{extarrows}
\usepackage{pifont}
\setlength{\topmargin}{0pt}
\setlength{\textheight}{9in}
\setlength{\headheight}{0pt}
\setlength{\oddsidemargin}{0.25in}
\setlength{\textwidth}{6in}
%%%%%%%%%%%%%%%%%%%%%%%%
%%%%%% Define math operator %%%%%
%%%%%%%%%%%%%%%%%%%%%%%%
\DeclareMathOperator*{\argmin}{\bf argmin}
\DeclareMathOperator*{\relint}{\bf relint\,}
\DeclareMathOperator*{\dom}{\bf dom\,}
\DeclareMathOperator*{\intp}{\bf int\,}
\DeclareMathOperator*{\aff}{\bf aff\,}
%%%%%%%%%%%%%%%%%%%%%%%%
\setlength{\topmargin}{0pt}
\setlength{\textheight}{9in}
\setlength{\headheight}{0pt}
\setlength{\oddsidemargin}{0.25in}
\setlength{\textwidth}{6in}
\pagestyle{fancy}
%%%%%%%%%%%%%%%%%%%%%%%%
%% Define the Exercise environment %%
%%%%%%%%%%%%%%%%%%%%%%%%
\mdtheorem[
topline=false,
rightline=false,
leftline=false,
bottomline=false,
leftmargin=-10,
rightmargin=-10
]{exercise}{\textbf{Exercise}}
%%%%%%%%%%%%%%%%%%%%%%%%
%% Define the Problem environment %%
%%%%%%%%%%%%%%%%%%%%%%%%
\mdtheorem[
topline=false,
rightline=false,
leftline=false,
bottomline=false,
leftmargin=-10,
rightmargin=-10
]{problem}{\textbf{Problem}}
%%%%%%%%%%%%%%%%%%%%%%%%
%% Define the Solution Environment %%
%%%%%%%%%%%%%%%%%%%%%%%%
\declaretheoremstyle
[
spaceabove=0pt,
spacebelow=0pt,
headfont=\normalfont\bfseries,
notefont=\mdseries,
notebraces={(}{)},
headpunct={:\;},
headindent={},
postheadspace={ },
bodyfont=\normalfont,
qed=$\blacksquare$,
preheadhook={\begin{mdframed}[style=myframedstyle]},
postfoothook=\end{mdframed},
]{mystyle}
\declaretheorem[style=mystyle,title=Solution]{solution}
\mdfdefinestyle{myframedstyle}{%
	topline=false,
	rightline=false,
	leftline=false,
	bottomline=false,
	skipabove=-6ex,
	leftmargin=-10,
	rightmargin=-10,
	backgroundcolor=black!5}
%%%%%%%%%%%%%%%%%%%%%%%%

\renewcommand{\eqref}[1]{Eq.~(\ref{#1})}
\newcommand{\figref}[1]{Fig.~\ref{#1}}
\newcommand{\proj}[2]{\textbf{P}_{#2} (#1)}
\newcommand{\lspan}[1]{\textbf{span}  (#1)  }
\newcommand{\rank}[1]{ \textbf{rank}  (#1)  }
\newcommand{\tr}{ \operatorname{tr}  }
\newcommand{\RNum}[1]{\uppercase\expandafter{\romannumeral #1\relax}}
\newcommand{\mb}[1]{\mathbf{#1}}
\DeclareMathOperator*{\cl}{\bf cl\,}
\DeclareMathOperator*{\bd}{\bf bd\,}
\DeclareMathOperator*{\conv}{\bf conv\,}
\DeclareMathOperator*{\epi}{\bf epi\,}
\theoremstyle{definition}	
\newtheorem{definition}{Definition}	

%%%%%%%%%%%%%%%%%%%%%%%%%%%%%%%%%%%%%%%%%%
%%           Homework info.             %%
\newcommand{\posted}{\text{Nov. 28th, 2025}}
\newcommand{\due}{\text{Dec. 16th, 2025}}
\newcommand{\hwno}{\text{3}}
%%%%%%%%%%%%%%%%%%%%%%%%%%%%%%%%%%%%%%%%%%

%%%%%%%%%%%%%%%%%%%%%%%%%%%%%%%%%%%%%
%%    Put your information here   %%
\newcommand{\name}{\text{Jinghao Chen}}
\newcommand{\id}{\text{PB23010359}}
%%%%%%%%%%%%%%%%%%%%%%%%%%%%%%%%%%%%

\chead{\textbf{Homework \hwno}}

\begin{document}
	\vspace*{-4\baselineskip}
	\thispagestyle{empty}
	
	\begin{center}
		{\bf\large Introduction to Machine Learning}\\
		{Fall 2025}\\
		University of Science and Technology of China
	\end{center}
	
	\noindent
	Lecturer: Zhihui Li, Xiaojun Chang
	\hfill
	Homework \hwno
	\\
	Posted: \posted
	\hfill
	Due: \due
    \\
    Name: \name
	\hfill
	ID: \id
	
	\noindent
	\rule{\textwidth}{2pt}
	
	\medskip
	
	\textbf{Notice, }to get the full credits, please present your solutions step by step.

    \begin{exercise}[Affine Sets and Projections]
        \begin{enumerate}[label=(\arabic*)]
        \item Let \(A\in\mathbb{R}^{m\times n}\), \(b\in\mathbb{R}^{m}\), and consider
        \[U=\{x\in\mathbb{R}^{n}:Ax=b\}.\]
        \begin{enumerate}[label=(\alph*)]
            \item Show that \(U\) is an affine set. Identify its direction space in terms of \(A\).
            \item Prove that for any \(x_{0}\in U\),
            \[U=x_{0}+\mathcal{N}(A),\]
            where \(\mathcal{N}(A)\) denotes the null space of \(A\).
        \end{enumerate}
        
        \item Suppose \(A\in\mathbb{R}^{m\times n}\) has full row rank and \(U=\{x:Ax=b\}\) is nonempty.
        \begin{enumerate}[label=(\alph*)]
            \item Consider the projection of a point \(y\in\mathbb{R}^{n}\) onto \(U\):
            \[p=\argmin_{x\in U}\|x-y\|^{2}_{2}.\]
            Use Lagrange multipliers to derive the formula
            \[p=y-A^{\top}(AA^{\top})^{-1}(Ay-b).\]
            \item In \(\mathbb{R}^{3}\), let
            \[U=\{x\in\mathbb{R}^{3}:x_{1}+2x_{2}-x_{3}=3\},\quad y=(1,0,0)^{\top}.\]
            Compute explicitly the projection \(p\) of \(y\) onto \(U\).
        \end{enumerate}
        \end{enumerate}
    \end{exercise}

\newpage

    \begin{solution}\textbf{Affine Sets and Projections}
        \begin{enumerate}[label=(\arabic*)]
        \item 
        \begin{enumerate}[label=(\alph*)]
            \item 
            $\forall$ $x_1, x_2 \in U$, $\forall \theta \in \mathbb{R}$, 
            \[
            A(\theta x_1 + (1-\theta)x_2) = \theta Ax_1 + (1-\theta)Ax_2 = \theta b + (1-\theta)b = b
            \Longrightarrow
            \theta x_1 + (1-\theta)x_2 \in U.
            \]
            namely $U$ is an affine set. 

            $\forall$ $x \in U$, $\exists$ $y \in \mathbb{R}^n$, s.t.
            \[
            x+y\in U
            \Longleftrightarrow
            A(x+y) = Ax+Ay = b+Ay = b
            \Longleftrightarrow
            Ay = 0
            \Longleftrightarrow
            y \in \mathcal{N}(A)
            \]
            namely the direction space of $U$ is $\mathcal{N}(A)$.
            \item 
            $\forall$ $x_0 \in U$,
            \[
            x\in U
            \Longleftrightarrow
            Ax = b
            \Longleftrightarrow
            A(x - x_0) = 0
            \Longleftrightarrow
            x - x_0 \in \mathcal{N}(A)
            \Longleftrightarrow
            x \in x_0 + \mathcal{N}(A)
            \]
            Thus, $U = x_0 + \mathcal{N}(A)$.
        \end{enumerate}
        \item 
        \begin{enumerate}[label=(\alph*)]
            \item 
            \begin{align*}
                \min_{x \in \mathbb{R}^n} \quad & \|x - y\|_2^2 \\
                \text{s.t.} \quad & A x = b
            \end{align*}
            \[
            \xRightarrow{def}
            L(x, \lambda) = \|x - y\|_2^2 + \lambda^{\top}(A x - b)
            \]
            \[
            \nabla_{x} L
            =
            2(x-y) + A^{\top} \lambda
            \xRightarrow{\nabla_{x} L = 0}
            p = x^* = y - \frac{1}{2} A^{\top} \lambda
            \]
            \[
            Ax = b
            \Longrightarrow
            A(y - \frac{1}{2} A^{\top} \lambda) = b
            \Longrightarrow
            A A^{\top} \lambda = 2(Ay - b)
            \]
            $A$ has full row rank $\Longrightarrow$ $AA^{\top} \in \mathbb{R}^{m \times m}$ is symmetric positive definite. $\Longrightarrow$ $AA^{\top}$ is invertible.
            \[
            \Longrightarrow
            \lambda
            =
            2(AA^{\top})^{-1}(Ay-b)
            \Longrightarrow
            p=y-A^{\top}(AA^{\top})^{-1}(Ay-b)
            \]
            \item
            \[
            A = (1,2,-1),\quad b=3,\quad y=(1,0,0)^{\top}.
            \]
            \[
            AA^{\top} = 1^2 + 2^2 + (-1)^2 = 6,\quad
            Ay - b = 1*1 + 2*0 + (-1)*0 - 3 = -2.
            \]
            \[
            p = y - A^{\top}(AA^{\top})^{-1}(Ay-b) 
            = 
            \begin{pmatrix}
                1\\
                0\\
                0
            \end{pmatrix} 
            - 
            \begin{pmatrix}
                1\\
                2\\
                -1
            \end{pmatrix}
            \cdot \frac{1}{6} \cdot (-2)
            = 
            \begin{pmatrix}
                \frac{4}{3}\\
                \frac{2}{3}\\
                -\frac{1}{3}
            \end{pmatrix}.
            \]
        \end{enumerate}
        \end{enumerate}
        \end{solution}

    \newpage

    \begin{exercise}[Convex Sets, Images and Preimages]
        Let \(C\subseteq\mathbb{R}^{n}\) be a nonempty convex set and let \(A\in\mathbb{R}^{m\times n}\), \(a\in\mathbb{R}^{m}\).
        \begin{enumerate}[label=(\arabic*)]
        \item Show that the affine image of \(C\),
        \[D=\{y\in\mathbb{R}^{m}:y=Ax+a,\,x\in C\},\]
        is convex.
        
        \item Define the affine preimage of \(C\) by
        \[E=\{y\in\mathbb{R}^{m}:Ay+a\in C\}.\]
        Show that \(E\) is convex.
        
        \item (Epigraph and hypograph)
        \begin{enumerate}[label=(\alph*)]
            \item Let \(f:\mathbb{R}^{n}\to\mathbb{R}\) be a convex function. The epigraph of \(f\) is 
            \[\epi f=\{(x,t)\in\mathbb{R}^{n+1}:t\geq f(x)\}.\]
            Show that \(\epi f\) is a convex set.
            \item Let \(g:\mathbb{R}^{n}\to\mathbb{R}\) be concave. Its hypograph is 
            \[\operatorname{hypo}g=\{(x,t)\in\mathbb{R}^{n+1}:t\leq g(x)\}.\]
            Show that \(\operatorname{hypo}g\) is convex.
        \end{enumerate}
        
        \item (Interior and relative interior -- computation) You do not need to prove your answers, but you should explain them briefly.
        \begin{enumerate}[label=(\alph*)]
            \item Find \(\intp C\) and \(\relint C\) for 
            \[C=\{x\in\mathbb{R}^{3}:x_{1}^{2}+x_{2}^{2}<4,\ x_{3}\geq 0\}.\]
            \item Let 
            \[K=\{x\in\mathbb{R}^{n}:x_{i}\geq 0,\ \sum_{i=1}^{n}x_{i}=1\}\]
            be the standard probability simplex. Determine \(\intp K\) (as a subset of \(\mathbb{R}^{n}\)) and \(\relint K\).
        \end{enumerate}
        \end{enumerate}
    \end{exercise}

\newpage

    \begin{solution}\textbf{Convex Sets, Images and Preimages}
        \begin{enumerate}[label=(\arabic*)]
            \item 
            $\forall$ $y_1, y_2 \in D$, $\exists$ $x_1, x_2 \in C$, s.t. $y_1 = A x_1 + a$, $y_2 = A x_2 + a$.
            \[
            C \text{ is convex } \Longrightarrow \forall \theta \in [0,1], \theta x_1 + (1-\theta) x_2 \in C,
            \]
            \[
            \Longrightarrow
            \theta y_1 + (1-\theta) y_2
            =
            \theta (A x_1 + a) + (1-\theta)(A x_2 + a)
            =
            A(\theta x_1 + (1-\theta) x_2) + a
            \in D
            \]
            namely $D$ is convex.
            \item 
            $\forall$ $y_1, y_2 \in E$, $Ay_1 + a \in C$, $Ay_2 + a \in C$,
            \[
            C \text{ is convex.}
            \Longrightarrow
            \forall \theta \in [0,1],
            \theta (Ay_1 + a) + (1-\theta)(Ay_2 + a)
            =
            A(\theta y_1 + (1-\theta) y_2) + a
            \in C
            \]
            $\Longrightarrow$ $\theta y_1 + (1-\theta) y_2 \in E$, namely $E$ is convex.
            \item 
            \begin{enumerate}[label=(\alph*)]
                \item 
                $\forall$ $(x_1, t_1), (x_2, t_2) \in \epi f$, $t_1 \geq f(x_1)$, $t_2 \geq f(x_2)$,
                \begin{align}
                f \text{ is convex }
                \Longrightarrow
                \forall \theta \in [0,1],
                \theta t_1 + (1-\theta) t_2 
                &\geq \theta f(x_1) + (1-\theta) f(x_2)\nonumber\\
                &\geq f(\theta x_1 + (1-\theta) x_2)\nonumber
                \end{align}
                $\Longrightarrow$ $(\theta x_1 + (1-\theta) x_2, \theta t_1 + (1-\theta) t_2) \in \epi f$, namely $\epi f$ is convex.
                \item
                \ding{172} $\forall$ $(x_1, t_1), (x_2, t_2) \in \operatorname{hypo} g$, $t_1 \leq g(x_1)$, $t_2 \leq g(x_2)$,
                \begin{align}
                g \text{ is concave }
                \Longrightarrow
                \forall \theta \in [0,1],
                \theta t_1 + (1-\theta) t_2 
                &\leq \theta g(x_1) + (1-\theta) g(x_2)\nonumber\\
                &\leq g(\theta x_1 + (1-\theta) x_2)\nonumber
                \end{align}
                $\Longrightarrow$ $(\theta x_1 + (1-\theta) x_2, \theta t_1 + (1-\theta) t_2) \in \operatorname{hypo} g$, namely $\operatorname{hypo} g$ is convex.

                \ding{173} Let $h = -g$, then $h$ is convex.
                \[
                t \leq g(x)
                \Longleftrightarrow
                -t \geq -g(x) = h(x)
                \]
                \[
                (x,t) \in \operatorname{hypo} g
                \Longleftrightarrow
                (x,-t) \in \operatorname{epi} h
                \]
                Define the linear mapping:
                \[
                T: \mathbb{R}^{n+1} \to \mathbb{R}^{n+1},\quad 
                T
                \begin{pmatrix}
                    x\\
                    t
                \end{pmatrix}
                =
                \begin{pmatrix}
                    I_n & 0\\
                    0 & -1
                \end{pmatrix} 
                \begin{pmatrix}
                    x\\
                    t
                \end{pmatrix}
                =
                \begin{pmatrix}
                    x\\
                    -t
                \end{pmatrix}
                \]
                According to (1) and (3)(a), $\epi h$ is convex $\Longrightarrow$ $T(\operatorname{hypo} g) = \epi h$ is convex $\Longrightarrow$ $\operatorname{hypo} g$ is convex.
            \end{enumerate}
            \item
            \begin{enumerate}[label=(\alph*)]
                \item 
                \[
                \intp C 
                =
                \relint C
                = 
                \{x\in\mathbb{R}^{3}:x_{1}^{2}+x_{2}^{2}<4,\ x_{3}>0\},
                \]
                Proof:
                \[
                B((0,0,1),\frac{1}{2}) \subset C
                \Longrightarrow
                3 \geq \dim C \geq \dim B((0,0,1),\frac{1}{2}) =3
                \]
                \[
                \Longrightarrow
                \dim C = \dim (\aff C) = 3
                \Longrightarrow
                \intp C = \relint C
                \]
                \item
                \[
                \intp K = \emptyset, \quad \relint K = \{x \in \mathbb{R}^n : x_i > 0, \sum_{i=1}^n x_i = 1\}.
                \]
                Proof:
                \[
                \dim K = n-1 < n
                \Longrightarrow
                \intp K = \emptyset
                \]
                \[
                \aff K = \{x \in \mathbb{R}^n : \sum_{i=1}^n x_i = 1\}
                \Longrightarrow
                \relint K = \{x \in \mathbb{R}^n : x_i > 0, \sum_{i=1}^n x_i = 1\}.
                \]
            \end{enumerate}
        \end{enumerate}
    \end{solution}

    \newpage

    \begin{exercise}[Cones and Farkas-type Results]
        Let \(A=(a_{1},\ldots,a_{n})\in\mathbb{R}^{m\times n}\) and consider the conic hull
        \[\operatorname{cone}(A)=\left\{\sum_{i=1}^{n}\alpha_{i}a_{i}:\alpha_{i}\geq 0\right\}\subseteq\mathbb{R}^{m}.\]
        
        \begin{enumerate}[label=(\arabic*)]
        \item Show that \(\operatorname{cone}(A)\) is a convex cone, i.e.,
        \[x,y\in\operatorname{cone}(A),\ \lambda,\mu\geq 0\ \Rightarrow\ \lambda x+\mu y\in\operatorname{cone}(A).\]
        
        \item Show that \(\operatorname{cone}(A)\) is closed.
        
        \item (Gordan's theorem) Let \(A\in\mathbb{R}^{m\times n}\). Prove that exactly one of the following two systems has a solution:
        \begin{itemize}
            \item There exists \(x\in\mathbb{R}^{n}\) such that \(Ax>0\) (componentwise).
            \item There exists \(y\in\mathbb{R}^{m}\), \(y\geq 0\), \(y\neq 0\), such that \(A^{\top}y=0\).
        \end{itemize}
        Hint: Consider appropriate cones and use separation theorems (you may assume standard separation results for closed convex cones).
        
        \item Show how Farkas' Lemma can be deduced from Gordan's theorem in part (3).
        \end{enumerate}
    \end{exercise}

\newpage

    \begin{solution}\textbf{Cones and Farkas-type Results}
        \begin{enumerate}[label=(\arabic*)]
            \item 
            $\forall$ $x,y \in \operatorname{cone}(A)$, $
            \exists$ $\alpha_i, \beta_i \geq 0$, s.t. $x = \sum_{i=1}^n \alpha_i a_i, \quad y = \sum_{i=1}^n \beta_i a_i$.
            \begin{align*}
            \forall \lambda, \mu \geq 0,
            \lambda \alpha_i + \mu \beta_i \geq 0,
            \Longrightarrow
            \lambda x + \mu y
            &=
            \lambda \sum_{i=1}^n \alpha_i a_i + \mu \sum_{i=1}^n \beta_i a_i\\
            &=
            \sum_{i=1}^n (\lambda \alpha_i + \mu \beta_i) a_i \in \operatorname{cone}(A)
            \end{align*}
            \item 
            First, prove that for each linear independent index set $I \subseteq \{1,2,\ldots,n\}$, the conic hull $\operatorname{cone}(A_I)$ is closed, where $A_I = (a_i)_{i \in I}$.

            $\forall$ $\{x_k\} \subset \operatorname{cone}(A_I)$ with $x_k \to x \in \mathbb{R}^m$.
            \[
            x_k = \sum_{i \in I} \alpha_{i}^{(k)} a_i, \quad \alpha_{i}^{(k)} \geq 0.
            \]
            If $\{\alpha_{i}^{(k)}\}$ is unbounded, then $\exists$ subsequence $\{\alpha_{i}^{(k_j)}\}$, s.t. $\left\|\alpha^{(k_j)}\right\| \to \infty$. Define the normalized sequence: $\beta^{(k_j)} = \frac{\alpha^{(k_j)}}{\left\|\alpha^{(k_j)}\right\|}$. Since $\beta^{(k_j)}$ is bounded, $\exists$ convergent subsequence $\beta^{(k_{j_l})} \to \beta \geq 0$, $\|\beta\| = 1$. According to the continuity of linear mapping, $A_I \beta^{(k_{j_l})} \to A_I \beta$. On the other hand, $A_I \beta^{(k_{j_l})} = \frac{x_{k_{j_l}}}{\left\|\alpha^{(k_{j_l})}\right\|} \to 0$. Thus, $A_I \beta = 0$, i.e., $\sum_{i \in I} \beta_i a_i = 0$, which contradicts the linear independence of $\{a_i\}_{i \in I}$.
            
            Therefore, $\{\alpha_{i}^{(k)}\}$ is bounded $\Longrightarrow$ $\exists$ convergent subsequence $\alpha^{(k_j)} \to \alpha \geq 0$. According to the continuity of linear mapping, $\lim_{j \to \infty} A_I \alpha^{(k_j)} = A_I \alpha \in \operatorname{cone}(A_I)$ $\Longrightarrow$ $x \in \operatorname{cone}(A_I)$ $\Longrightarrow$ $\operatorname{cone}(A_I)$ is closed.

            Second, every point in $C$ can be expressed as a nonnegative combination of at most $m$ linearly independent columns of $A$. Therefore, there are only finite linear independent index sets. For each linear independent index set $I$, $\operatorname{cone}(A_I)$ is closed $\Longrightarrow$ $\operatorname{cone}(A) = \bigcup_{I} \operatorname{cone}(A_I)$ is closed.
            \item 
            On the one hand, assume the first one has no solution. Then the sets
            \[
            U = \{Ax: x\in\mathbb{R}^n\}
            , \quad
            V = \{y \in \mathbb{R}^m: y > 0\}
            \]
            are disjoint convex sets. By the separation theorem for convex sets, $\exists$ $a \neq 0 \in \mathbb{R}^m$ and $\alpha \in \mathbb{R}$, s.t.
            \[
            a^{\top} u \leq \alpha, \forall u \in U
            ,\quad
            a^{\top} v \geq \alpha, \forall v \in V
            \]
            $0 \in U$, $a^{\top} 0 = 0 \leq \alpha$ $\Longrightarrow$ $\alpha \geq 0$.
            $\forall$ $v \in V$, let $v \to 0$, then $\alpha \leq a^{\top} v \to 0$ $\Longrightarrow$ $\alpha = 0$.

            Besides, 
            \[
            \text{``}u \in U
            \Longleftrightarrow
            -u \in U\text{''}
            \Longrightarrow
            a^{\top} u \leq 0, a^{\top} (-u) \leq 0
            \Longrightarrow
            a^{\top} u = 0, \forall u \in U
            \]
            \[
            \left.
            \begin{aligned}
            a^{\top} u = a^{\top} A x = 0, \forall x\in \mathbb{R}^n &\Longrightarrow A^{\top} a = 0\\
            a^{\top} v \geq 0, \forall v \in V &\Longrightarrow a \geq 0
            \end{aligned}
            \right\}
            \Longrightarrow
            \exists y = a \geq 0, y \neq 0, \text{ s.t. } A^{\top} y = 0.
            \]
            On the other hand, assume the second one holds. Then 
            \[
            \exists y \in \mathbb{R}^m, y\neq 0, y \geq 0, \text{ s.t. } A^{\top} y = 0.
            \Longrightarrow 
            \forall x \in \mathbb{R}^n, \quad y^{\top} A x = x^{\top} A^{\top} y = 0
            \]
            Because $y \neq 0$ and $y \geq 0$, $\exists$ $i$, s.t. $y_i > 0$ $\Longrightarrow$ $A x > 0$ is impossible. Thus, the first one has no solution.

            In conclusion, the first one and the second one are mutually exclusive and collectively exhaustive, i.e., exactly one of the two systems has a solution.
            \item
            Farkas' Lemma:
            Let $A \in \mathbb{R}^{m \times n}$ and $b \in \mathbb{R}^m$. Then exactly one of the following two systems has a solution:
            \begin{itemize}
                \item There exists $x \in \mathbb{R}^n$ such that $Ax = b$, $x \geq 0$ (componentwise).
                \item There exists $y \in \mathbb{R}^m$ such that $A^{\top} y \geq 0$, $b^{\top} y < 0$.
            \end{itemize}
            Deduce Farkas' Lemma from Gordan's theorem:

            Construct the augmented matrix $\tilde{A} = [A, -b]^{\top} \in \mathbb{R}^{(n+1) \times m}$. 

            Case 1: $\exists$ $y \in \mathbb{R}^m$, s.t. $\tilde{A} y > 0$ $\Longleftrightarrow$ $A^{\top} y > 0$, $-b^{\top} y > 0$ $\Longleftrightarrow$ $A^{\top} y > 0$, $b^{\top} y < 0$ $\Longrightarrow$ the second one in Farkas' Lemma has a solution.

            Case 2: $\exists$ $v \in \mathbb{R}^{n+1}$, $v \geq 0$, $v \neq 0$, s.t. $\tilde{A}^{\top} v = 0$. Define $v = (x,t)$, $x \in \mathbb{R}^n$, $t \in \mathbb{R}$. Then
            \[
            \tilde{A} v
            =
            A x - b t
            =
            0
            ,
            (x,t) \geq 0, (x,t) \neq 0
            \]
            If $t > 0$, let $x' = \frac{x}{t}$, then
            \[
            A x' = b, \quad x' \geq 0,
            \]
            namely the first one in Farkas' Lemma has a solution.

            If $t = 0$, then $A x = 0$, $x \geq 0$, $x \neq 0$. If the first one has solution, it does not matter. If the first one has no solution, it means $b \notin \operatorname{cone}(A)$. According to (1) and (2), $\operatorname{cone}(A)$ is convex and closed. By the separation theorem for closed convex sets, $\exists$ $y \neq 0 \in \mathbb{R}^m$ and $\alpha \in \mathbb{R}$, s.t.
            \[
            y^{\top} u \geq \alpha, \forall u \in \operatorname{cone}(A)
            ,\quad
            y^{\top} b < \alpha
            \]
            $0 \in \operatorname{cone}(A)$ $\Longrightarrow$ $\alpha \leq 0$. If $\exists$ $u_0\in \operatorname{cone}(A)$, s.t. $y^{\top} u_0 < 0$, then for $\forall$ $\lambda > 0$, $\lambda u_0 \in \operatorname{cone}(A)$, $y^{\top} (\lambda u_0) = \lambda (y^{\top} u_0) \to -\infty$, which contradicts $y^{\top} u \geq \alpha$. Thus, 
            \[
            y^{\top} u \geq 0, \forall u \in \operatorname{cone}(A)
            ,\quad
            y^{\top} b < \alpha \leq 0
            \]
            \[
            a_i = A e_i \in \operatorname{cone}(A)
            \Longrightarrow
            y^{\top} a_i \geq 0, \forall i
            \Longrightarrow
            A^{\top} y \geq 0
            ,\quad
            b^{\top} y = y^{\top} b < 0
            \]
            namely the second one in Farkas' Lemma has a solution.
        \end{enumerate}
    \end{solution}

    \newpage
    \bibliography{refs}
    \bibliographystyle{abbrv}
	
\end{document}